La supervisión débil mejora la IoU principal en 0 de tres pérdidas en ID y en 2 de tres en Bolivia. La matriz permite observar sensibilidad a la pérdida y a la fuente de etiquetas, sin demostrar una mejora universal ni aislar el efecto de las etiquetas del número de actualizaciones. Los casos cualitativos muestran ambigüedad radiométrica, omisión de agua de respuesta alta y errores de detalle espacial. La evaluación externa debe conservar la agregación y el signo real de la brecha.

